% Editable reconstruction of the existing CV PDF, revised on 2026-10-15.
% This is a new source file, not the recovered original TeX.
% Original PDF: Woochang_CV.pdf (preserved), SHA-256:
% df3b6429d030cc74fc616445ba3db5fedf01ae32294d91eac7c50a7334d27cb6
% Completed-work facts, dates, skills, and metrics come from that PDF.
% Current-work scope comes from the Harness and Jev/Kev chats and their
% project records; no unpublished performance gains are claimed here.
\documentclass[10pt,a4paper]{article}
\usepackage[T1]{fontenc}
\usepackage[utf8]{inputenc}
\usepackage{lmodern}
\usepackage[margin=14mm]{geometry}
\usepackage{enumitem,titlesec,microtype,xcolor}
\usepackage{hyperref}
\definecolor{linkblue}{HTML}{17467D}
\hypersetup{colorlinks=true,urlcolor=linkblue,linkcolor=linkblue,
  pdftitle={Woochang Sim - Research CV},pdfauthor={Woochang Sim},
  pdfsubject={LLM reasoning evaluation, agent harnesses, and self-monitoring}}
\input{glyphtounicode}
\pdfgentounicode=1
\urlstyle{same}
\pagestyle{empty}
\setlength{\parindent}{0pt}
\setlength{\parskip}{3pt}
\setlength{\emergencystretch}{2em}
\titleformat{\section}{\large\scshape}{}{0pt}{}[\titlerule]
\titlespacing*{\section}{0pt}{8pt}{4pt}
\setlist[itemize]{leftmargin=12pt,label=\textendash,
  itemsep=1pt,topsep=2pt,parsep=0pt,partopsep=0pt}
\newcommand{\entry}[3]{\textbf{#1}\hfill\textit{#2}\\\textit{#3}\par}
\newcommand{\pub}[2]{\noindent #1\\{\small #2}\par\vspace{2pt}}

\begin{document}
\begin{center}
  {\LARGE\bfseries Woochang Sim}\\[3pt]
  Ph.D. Student in AI Convergence, GIST Data Science Lab\\[3pt]
  \href{https://whatchang.github.io/}{whatchang.github.io}
  \enspace|\enspace\href{mailto:woochang@gm.gist.ac.kr}{woochang@gm.gist.ac.kr}\\[2pt]
  {\small\href{https://github.com/whatchang}{github.com/whatchang}
  \enspace|\enspace\href{https://www.linkedin.com/in/woo-chang-sim-b739a8143/}{linkedin.com/in/woo-chang-sim-b739a8143}}
\end{center}
\vspace{-5pt}

\section{Research Profile}
Ph.D. researcher with published work in \textbf{controlled LLM reasoning evaluation}
and experience developing \textbf{executable synthetic benchmarks}.
My current research studies diagnostics for automated LLM-agent harness improvement
and learned internal readouts for model judgment and self-monitoring.
I connect these questions through algorithm implementation, controlled experiments,
and analysis of model behavior; adaptive use of computation is a further research goal.

{\small\textbf{Broader interests:} Adaptive agents under uncertainty and distribution shift;
probabilistic task and environment modeling; Bayesian and meta-reinforcement learning;
generalization to unseen tasks and environments; language-model priors and world models
for interactive adaptation.}

\section{Education}
\textbf{Ph.D. in AI Convergence}, Gwangju Institute of Science and Technology\hfill 2025--Present\\
Advisor: Prof. Sundong Kim\par
\textbf{M.S. in Artificial Intelligence}, Gwangju Institute of Science and Technology\hfill 2023--2025\\
Advisor: Prof. Sundong Kim\par
\textbf{B.S. in Computer Engineering}, Chosun University\hfill 2016--2023

\section{Research Experience}
\entry{Harness Diagnostics for Interactive LLM Agents}{2026--Present}{Ongoing Ph.D. research; ARC-AGI-3 testbed}
\begin{itemize}
  \item Developed and piloted diagnostic tooling that exposes solver token probabilities
  and traces recorded context-production paths to support automated harness edits.
  \item Designed controlled comparisons that separate diagnostic judgments from executed
  behavioral outcomes and keep development tasks separate from final evaluation tasks.
  \item Analyzed recurring agent failures and incomplete provenance mappings;
  tool use and repair effects remain under pilot evaluation.
\end{itemize}

\entry{Internal Readouts for Model Judgment and Self-Monitoring}{2026--Present}{Ongoing Ph.D. research}
\begin{itemize}
  \item Investigating learned readouts from unused output coordinates for model judgment
  while preserving language-generation ability.
  \item Studying whether these signals can track reasoning progress and inform self-checking
  and future compute allocation; signal validity, intervention effects, and allocation
  benefits are separate research questions.
\end{itemize}

\entry{Controlled Evaluation of LLM Reasoning on ARC}{Nov 2023--Jan 2025}{Co-first author, ACM Transactions on Intelligent Systems and Technology (TIST)}
\begin{itemize}
  \item Operationalized a Language of Thought framework into reproducible experiments
  on \textbf{logical coherence, compositionality, and productivity}.
  \item Led implementation of controlled reasoning experiments, including prompting-based
  evaluations, an action-based compositionality environment, and an interface for human baselines.
  \item Compared language models with human performance and analyzed gaps between
  task accuracy and valid abstract and compositional reasoning.
\end{itemize}

\entry{GIFARC: Synthetic Analogy-Based ARC Tasks}{Apr 2025--Present}{arXiv preprint; \href{https://gifarc.vercel.app/}{project}}
\begin{itemize}
  \item Developed a \textbf{three-stage VLM/LLM generation pipeline} that converts GIF-derived
  analogies into visual abstractions, task sketches, and executable ARC-style task instances.
  \item Evaluated analogical information through controlled \textbf{in-context learning and
  supervised fine-tuning} experiments, including downstream ARC benchmarks.
  \item Coordinated experimental planning, task allocation, and weekly research
  discussions across collaborators.
\end{itemize}

\newpage
\section{Publications}
{\small $*$ denotes equal contribution.\par}
\pub{\textbf{Reasoning Abilities of Large Language Models: In-Depth Analysis on the Abstraction and Reasoning Corpus}}
{S. Lee$^*$, \textbf{W. Sim}$^*$, D. Shin$^*$, W. Seo, J. Park, S. Lee, S. Hwang, S. Kim, S. Kim.\\
ACM TIST, 2025. \href{https://doi.org/10.1145/3712701}{doi:10.1145/3712701};
\href{https://github.com/GIST-DSLab/ARC_Prompt}{code}.}

\pub{\textbf{GIFARC: Synthetic Dataset for Leveraging Human-Intuitive Analogies to Elevate AI Reasoning}}
{\textbf{W. Sim}$^*$, H. Ryu$^*$, K. Choi$^*$, S. Han, S. Kim.\\
arXiv preprint, 2025. \href{https://arxiv.org/abs/2505.20672}{arXiv:2505.20672}.}

\pub{\textbf{Reasoning Abilities of Large Language Models through the Lens of Abstraction and Reasoning}}
{S. Lee, \textbf{W. Sim}, D. Shin, S. Kim, S. Kim.\\
NeurIPS System-2 Reasoning at Scale \textbf{Workshop}, 2024.
\href{https://openreview.net/forum?id=getEDedhSV}{paper}.}

\pub{\textbf{Enhancing Analogical Reasoning in the Abstraction and Reasoning Corpus via Model-Based RL}}
{J. Lee$^*$, \textbf{W. Sim}$^*$, S. Kim, S. Kim.\\
IJCAI IARML \textbf{Workshop}, 2024. \href{https://arxiv.org/abs/2408.14855}{arXiv:2408.14855}.}

\pub{\textbf{The Possibility of Prompt Engineering for ARC Problem Solving}}
{\textbf{W. Sim}, H. Jin, S. Kim, S. Kim.\\
Journal of KIISE: Computing Practices and Letters, 30(2):63--69, 2024.
\href{https://www.dbpia.co.kr/journal/articleDetail?nodeId=NODE11716261}{paper}.}

\pub{\textbf{Using Contrastive Learning for Abstraction and Reasoning Task}}
{G. Gu, \textbf{W. Sim}, J. Im, S. Kim, S. Kim.\\
Korea Software Congress, pp. 828--830, 2023.
\href{https://www.dbpia.co.kr/journal/articleDetail?nodeId=NODE11705255}{paper}.}

\pub{\textbf{Augmenting Few-Shot Demonstrations with Large Language Model}}
{W. Seo, \textbf{W. Sim}, S. Kim.\\
Korea Software Congress, pp. 341--343, 2023.
\href{https://www.dbpia.co.kr/journal/articleDetail?nodeId=NODE11705093}{paper}.}

\section{Technical Skills}
\textbf{Programming:} Python.\\
\textbf{Frameworks and tools:} PyTorch, Hugging Face Transformers, Weights \& Biases, Git.\\
\textbf{Research methods:} Experimental design, benchmark construction, ablation studies,
statistical evaluation, human evaluation, inter-annotator agreement analysis.\\
\textbf{ML experience:} LLM fine-tuning, in-context learning, synthetic data generation,
information retrieval, open-domain question answering, data pipelines.\\
\textbf{Languages:} Korean (native), English (intermediate).

\section{Selected Projects}
\textbf{Open-Domain Question Answering System}
\href{https://github.com/whatchang/mrc-level2-nlp-04}{[code]}\hfill Oct--Nov 2021\\
Implemented a hybrid BM25/Elasticsearch/DPR retrieval pipeline and a RoBERTa-large
reader over 60K Wikipedia documents; achieved \textbf{68.89 EM / 78.69 F1},
ranking 6th among 19 teams.\par
\textbf{Relation Extraction with Typed Entity Markers}
\href{https://github.com/whatchang/klue-level2-nlp-04}{[code]}\hfill Sep--Oct 2021\\
Implemented typed entity markers and domain-adaptive masked-language-model
pretraining with RoBERTa-large; achieved \textbf{72.6 F1}.\par
\textbf{K-pop Relation Extraction Dataset Construction}\hfill Nov 2021\\
Constructed a 2,468-instance dataset across 10 relation classes and developed
iterative annotation guidelines; achieved Fleiss' $\kappa = 0.842$.

\section{Teaching and Service}
\textbf{Teaching Assistant, Machine Learning \& Deep Learning, GIST}\hfill Sep--Dec 2025\\
Graded machine-learning assignments with detailed feedback and taught fundamental
neural-network, convolutional-network, and recurrent-network concepts.\par
\textbf{Student Representative, AI Convergence Department, GIST}\hfill Jan--Dec 2025\\
Organized student events and facilitated communication between students and faculty.

\section{Awards}
The Chungshim Fellowship\hfill Mar 2024\\
Best Presentation Award, Korea Computer Congress 2023\hfill Jul 2023
\end{document}
